\section{Materials and Methods}
\label{sec:idp-materials-methods}

This section outlines the comprehensive experimental methodology employed to evaluate the performance and efficacy of memory-aware chunking.
The experiment is structured as a multi-stage pipeline that encompasses data generation, systematic computational profiling within isolated Docker containers, rigorous aggregation of profiling results, and detailed analysis supported by visualization.

\subsection*{Overall Structure}

The experimental pipeline is orchestrated by a single Bash script, \texttt{experiment.sh}, which sequentially coordinates four core Python scripts:

\begin{enumerate}
    \item \textbf{Data Generation} (\texttt{generate\_data.py})
    \item \textbf{Profiling Execution} (\texttt{collect\_profile.py})
    \item \textbf{Result Consolidation} (\texttt{collect\_results.py})
    \item \textbf{Analysis and Plotting} (\texttt{analyze\_results.py})
\end{enumerate}

Each step executes within \ac{DIND} containers, ensuring consistent environments, reproducibility of results, and clear isolation of software dependencies.

\subsection*{Configuration and Environment}

All experiments were executed on a single \ac{HPC} node at the \ac{UNICAMP} Discovery Labs, equipped with an Intel\textregistered\ Xeon\textregistered\ Silver 4310 processor, 256~\ac{GB} of \ac{RAM}, and two \ac{RTX}~A6000 \ac{GPU}s. Despite the presence of \ac{GPU} resources, all runs relied exclusively on \ac{CPU} execution to ensure comparability across input shapes.

The experiment is extensively configurable through environment variables specified at the beginning of the orchestrator script.
Important configurable parameters include:

\begin{itemize}
    \item \texttt{TIMESTAMP}: Unique identifier for output directories, defaulting to the current date and time.
    \item \texttt{CPUSET\_CPUS}: Specifies the CPU cores available to Docker, defaulting to core 0.
    \item \texttt{EXPERIMENT\_IMAGE\_TAG}: Docker image tag used for executing Python scripts.
    \item \texttt{EXPERIMENT\_N\_RUNS}: Number of repetitions for each scenario, defaulting to 3 runs per configuration.
    \item \texttt{DATASET\_INITIAL\_SIZE}, \texttt{DATASET\_FINAL\_SIZE}, \texttt{DATASET\_STEP\_SIZE}: Control the dimensions of synthetic seismic data generated, ranging from 100³ to 400³ voxels, incremented by steps of 100.
    \item \texttt{WORKER\_SCENARIOS}: Specifies worker configurations (\eg, \texttt{single:1}, \texttt{two:2}, \texttt{smalln:4}, \texttt{bign:8}).
    \item \texttt{CHUNKING\_MODES}: Determines which chunking strategies to compare, including \texttt{auto}, \texttt{evenly\_split}, and \texttt{memaware}.
    \item \texttt{GST3D\_MODEL\_FILE}: Path to the pre-trained predictive model for memory-aware chunking.
    \item \texttt{MEMORY\_LIMIT\_GB}: Defines the total memory allocation available per scenario, partitioned evenly among workers.
\end{itemize}

All configuration settings are logged at the experiment’s initiation for transparency and debugging purposes.

\subsection*{Step 1: Data Generation}

Data generation utilizes the \texttt{generate\_data.py} script, which systematically creates synthetic \ac{SEG-Y} seismic datasets across a grid of specified shapes.
It iteratively generates data volumes for all combinations of dimensions defined within the configured size range (\eg, $(100,100,100)$ through $(400,400,400)$), saving each as individual \texttt{.segy} files.

\subsection*{Step 2: Profiling Scenarios}

This stage represents the experiment’s core, systematically evaluating the computational performance of different chunking methods.
For each synthetic dataset, the pipeline:

\begin{enumerate}
    \item Deploys a local Dask cluster within a containerized environment, with a specified number of workers and memory allocation per worker derived from \texttt{MEMORY\_LIMIT\_GB}.
    \item Tests each chunking strategy:\begin{itemize}
                                           \item \texttt{auto}: Default heuristic-driven chunking provided by Dask.
                                           \item \texttt{evenly\_split}: Divides volumes into equal-sized sub-volumes.
                                           \item \texttt{memaware}: Employs the pre-trained memory prediction model to estimate the largest feasible chunk dimension that safely fits within memory constraints, adjusted by a safety factor.
    \end{itemize}
    \item Executes the gradient-structure-tensor computational kernel while concurrently recording execution time and detailed memory usage profiles (peak and average usage) via periodic sampling.
    \item Stores profiling results into structured JSON files, systematically named by dataset shape, chunking mode, worker count, and timestamp.
\end{enumerate}

\subsection*{Step 3: Result Consolidation}

The \texttt{collect\_results.py} script aggregates the profiling JSON outputs into comprehensive, structured CSV datasets:

\begin{itemize}
    \item \texttt{profiles\_summary.csv}: Summarizes each experimental run, including total execution time, peak memory usage, and \ac{OOM} occurrences.
    \item \texttt{profiles\_detail.csv}: Provides detailed, per-worker memory usage histories and averages for in-depth analysis.
\end{itemize}

This consolidation step also addresses incomplete experiments by marking missing scenario results explicitly as \ac{OOM} or failures, ensuring data completeness.

\subsection*{Step 4: Analysis and Visualization}

The final analytical step, executed by \texttt{analyze\_results.py}, interprets consolidated data to assess chunking strategies comprehensively:

\begin{itemize}
    \item Generates visualizations, including line charts comparing execution time and peak memory usage across worker counts and chunking strategies.
    \item Produces detailed box plots for per-worker memory usage, facilitating deeper insight into workload distribution.
    \item Creates a leaderboard \ac{CSV} summarizing the best-performing chunking mode for each dataset shape and worker scenario combination, quantifying relative performance improvements.
    \item Aggregates additional statistical summaries such as means, medians, \ac{OOM} frequencies, and distributions of chunk counts.
\end{itemize}

\subsection*{Experimental Significance}

The structured approach described above enables rigorous comparative analysis between traditional heuristic chunking and the proposed memory-aware strategy.
Key benefits include:

\begin{itemize}
    \item \textbf{Automated Reproducibility}: A single, comprehensive script facilitates repeatable, end-to-end experimentation.
    \item \textbf{Fair Comparison}: Identical datasets, hardware constraints, and repetitions ensure unbiased performance evaluations.
    \item \textbf{Memory-Awareness}: Empirically tests the practical benefits of leveraging predictive modeling to select optimal chunk sizes.
    \item \textbf{Comprehensive Metrics}: Captures detailed performance (execution time) and resource utilization (memory usage), enabling nuanced performance assessment.
\end{itemize}

In summary, this detailed experimental methodology enables robust validation of memory-aware chunking as a superior approach to managing computational resources effectively and enhancing data parallelism in high-performance computing contexts.